\DocumentMetadata{
  pdfversion=2.0,
  pdfstandard=ua-2,
  lang=en-US,
  tagging=on,
  tagging-setup={math/setup=mathml-SE,tabsorder=structure}
}
\documentclass[11pt,letterpaper]{article}
\usepackage[margin=0.68in,includefoot]{geometry}
\usepackage{newcomputermodern}
\usepackage{amsmath,amscd,mathtools}
\usepackage{tikz,adjustbox}
\providecommand{\square}{\mdlgwhtsquare}
\usepackage[hidelinks]{hyperref}
\hypersetup{
  pdftitle={Algebra First-Year Exam, May 2018, Part 2},
  pdfauthor={University of Florida Department of Mathematics},
  pdfsubject={Graduate mathematics examination},
  pdfdisplaydoctitle=true,
  bookmarksopen=true
}
\setcounter{secnumdepth}{0}
\setlength{\parindent}{0pt}
\setlength{\parskip}{0.28em}
\setlength{\emergencystretch}{3em}
\setlength{\leftmargini}{2.15em}
\setlength{\leftmarginii}{2.65em}
\setlength{\leftmarginiii}{3em}
\pagestyle{empty}
\raggedbottom
\allowdisplaybreaks
\NewTaggingSocketPlug{block/list/label}{examproblem}{%
  \tagstructbegin{tag=Lbl}\tagstructend%
  \tagstructbegin{tag=\LBody}%
  \tagstructbegin{tag=\ProblemHeadingTag,title-o={\ProblemHeadingText}}%
  \tagmcbegin{tag=\ProblemHeadingTag,actualtext={\ProblemHeadingText}}%
  #2%
  \tagmcend%
  \tagstructend%
}
\newenvironment{examproblems}[2]{%
  \def\ProblemHeadingTag{#2}%
  \AssignTaggingSocketPlug{block/list/label}{examproblem}%
  \begin{enumerate}%
  \setcounter{enumi}{#1}%
  \renewcommand{\labelenumi}{\textbf{\theenumi.}}%
  \setlength{\topsep}{0.35em}%
  \setlength{\partopsep}{0pt}%
  \setlength{\itemsep}{0.48em}%
  \setlength{\parsep}{0.18em}%
}{\end{enumerate}}
\newenvironment{examsubparts}{%
  \AssignTaggingSocketPlug{block/list/label}{default}%
  \begin{enumerate}%
  \setlength{\topsep}{0.18em}%
  \setlength{\partopsep}{0pt}%
  \setlength{\itemsep}{0.16em}%
  \setlength{\parsep}{0.1em}%
}{\end{enumerate}}
\newenvironment{examinstructions}{%
  \par\begingroup\itshape
}{%
  \par\endgroup\vspace{0.18em}
}
\makeatletter
\renewcommand\subsection{\@startsection{subsection}{2}{0pt}%
  {-1.25ex plus -.25ex minus -.1ex}%
  {0.55ex plus .1ex}%
  {\normalfont\large\bfseries}}
\renewcommand{\@maketitle}{%
  \newpage\null\vskip -1em%
  {\footnotesize\bfseries
    \href{https://www.ufl.edu/}{University of Florida}, \href{https://math.ufl.edu/}{Department of Mathematics}\par}%
  \vskip -0.5em%
  \tagstructbegin{tag=H1,title={Algebra First-Year Exam, May 2018, Part 2}}%
  \tagpdfparaOff%
  \tagmcbegin{tag=H1}%
    {\LARGE\bfseries\href{https://gma.math.ufl.edu/exams/algebra-first-year/}{Algebra First-Year Exam}, May 2018, Part 2\par}%
  \tagmcend%
  \tagpdfparaOn%
  \tagstructend%
  \vskip 0.25em%
  \tagmcbegin{artifact}%
    \hrule height 1pt%
  \tagmcend%
  \vskip 1em%
}
\makeatother
\title{Algebra First-Year Exam, May 2018, Part 2}
\author{}
\date{}
\begin{document}
\maketitle
\thispagestyle{empty}
\begin{examinstructions}
Answer four problems. Write your solutions clearly in complete English sentences. You may quote results (within reason) as long as you state them clearly.
\end{examinstructions}
\begin{examproblems}{0}{H2}
\def\ProblemHeadingText{Problem 1.}
\item
Let~\(D\) be a positive integer such that \(-D \equiv 1 \pmod 4\). Define
\[
\omega=\frac{1+\sqrt{-D}}{2}, \qquad R=\{a+b\omega:a,b\in\mathbb Z\}.
\]
 For \(\alpha\in R\) define \(\operatorname{N}(\alpha)=\alpha\overline{\alpha}\), where \(\overline{\alpha}\) is the complex conjugate of~\(\alpha\).
\begin{examsubparts}
\item[\textnormal{(a)}]
Prove that~\(R\) is a subring of the field \(\mathbb C\) of complex numbers.
\item[\textnormal{(b)}]
Prove that for all \(\alpha\in R\) we have \(\operatorname{N}(\alpha)\in\mathbb Z_{\ge 0}\).
\item[\textnormal{(c)}]
Prove that for all \(\alpha,\beta\in R\) we have \(\operatorname{N}(\alpha\beta)=\operatorname{N}(\alpha)\operatorname{N}(\beta)\).
\item[\textnormal{(d)}]
Let \(\alpha\in R\). Prove that~\(\alpha\) is a unit if and only if \(\operatorname{N}(\alpha)=1\).
\end{examsubparts}
\def\ProblemHeadingText{Problem 2.}
\item
This problem has two parts.
\begin{examsubparts}
\item[\textnormal{(a)}]
Let~\(F\) be a field and let~\(G\) be a finite subgroup of the unit group \(F^\times\). Prove that~\(G\) is cyclic.
\item[\textnormal{(b)}]
Give an example of a field~\(F\) and a finitely generated subgroup~\(G\) of \(F^\times\) which is not cyclic.
\end{examsubparts}
\def\ProblemHeadingText{Problem 3.}
\item
Let~\(R\) be a commutative ring with \(1\ne 0\) and let~\(M\) be a unital~\(R\)-module. Say~\(M\) is irreducible if \(M\ne\{0\}\) and the only submodules of~\(M\) are \(\{0\}\) and~\(M\). Prove that if~\(M\) is an irreducible~\(R\)-module then \(M\cong R/I\) for some maximal ideal~\(I\) in~\(R\).
\def\ProblemHeadingText{Problem 4.}
\item
Let~\(V\) be a vector space over a field~\(F\) and let~\(S\) be a subset of~\(V\) such that \(\operatorname{Span}(S)=V\). Use Zorn’s lemma to prove that there is a subset \(\mathcal B\) of~\(S\) which is a basis for~\(V\).
\def\ProblemHeadingText{Problem 5.}
\item
Find a representative for each similarity class of \(3\times 3\) matrices~\(A\) with entries in \(\mathbb F_3=\mathbb Z/3\mathbb Z\) such that \(A^4=A\). (Note that \(X^4-X\) factors into linear factors over \(\mathbb F_3\).)
\def\ProblemHeadingText{Problem 6.}
\item
Let \(\gamma\in\mathbb C\) be a root of \(X^3-6X^2+2\in\mathbb Q[X]\). Express \((1+\gamma^2)(3-4\gamma^2)\) and \(1/(1+\gamma^2)\) in the form \(a+b\gamma+c\gamma^2\), with \(a,b,c\in\mathbb Q\).
\end{examproblems}
\end{document}
